\chapter*{使用说明与真题总览}
\addcontentsline{toc}{chapter}{使用说明与真题总览}
\markboth{使用说明与真题总览}{}

\section*{一、本册怎么用}

\begin{enumerate}
  \item \textbf{第一遍（限时模考）}：按 150 分钟（3 小时）计时，从"历年试题"部分选一套完整做，
        答题时强迫自己写出每一步公式与单位——818 计算题占 58 分，过程分很重。
  \item \textbf{第二遍（对答案）}：翻到书末"参考答案"，逐题核对\textbf{公式选取、断面选取、单位}三个环节，
        不要只看最终数值。
  \item \textbf{第三遍（统计）}：把做错的题按考点归类，回到《教材精讲》对应的考点框重看"含义"与"适用条件"。
  \item \textbf{第四遍（考前一周）}：只重做近 5 年的选择、填空、判断（共 80 分），
        这部分是"真题重复率很高"的最大受益区。
\end{enumerate}

\section*{二、真题年份与题型一览}

\begin{center}
\begin{tabularx}{\textwidth}{@{}clX@{}}
\toprule
年份 & 科目代码 & 题型与分值 \\
\midrule
2014 & 818 & 选择 30（10$\times$3）+ 填空 30（8 题共 10 空$\times$3）+ 计算 90（1--7 必做，8--11 任选两道）；含少量明渠、流体机械内容 \\
2015 & 819 & 选择 45（15$\times$3）+ 填空 30（7 题共 10 空$\times$3）+ 计算 75（1--7 必做，8--10 任选一道）；含明渠内容 \\
2016 & 817 & 选择 30（10$\times$3）+ 填空 30（10 空$\times$3）+ 计算 90（1--7 必做，8--11 任选两道）；**本年无判断题** \\
2017 & 819 & 选择 30（10$\times$3）+ 填空 30（10 空$\times$3）+ 判断 30（10$\times$3）+ 计算 60（6 题） \\
2018 & 819 & 选择 15（5$\times$3）+ 判断 20（10$\times$2）+ 填空 30（8 题）+ 计算 85（1--7 必做 65，8--11 任选两道各 10）；**原卷判断题排在填空题之前** \\
2019 & 821 & 选择 10（5$\times$2）+ 判断 20（10$\times$2）+ 填空 30（6 题共 10 空$\times$3）+ \textbf{作图 7} + 简答 16（6+4+3+3）+ \textbf{证明 17（7+10）} + 计算 50（7 题各 10，1--4 必做，5--7 任选一道） \\
2020 & 818 & 选择 20（10$\times$2）+ 判断 10（10$\times$1）+ 填空 20（10 空$\times$2）+ 简答 20（6+6+8）+ \textbf{证明 20（2$\times$10）} + 计算 60（8/12/10/14/8/8） \\
2021 & 818 & 选择 30（10$\times$3）+ 判断 20（10$\times$2）+ 填空 30（10 空$\times$3）+ 简答 12（3$\times$4）+ 计算 58（10+8+14+8+10+8） \\
2022 & 818 & 选择 30（10$\times$3）+ 判断 20（10$\times$2）+ 填空 30（10 空$\times$3）+ 简答 12（3$\times$4）+ 计算 58（8/12/10/10/12/6） \\
2026 & 818 & 回忆版：选择、填空"与往年真题和题库类似"；简答 3 题；计算 5 题 \\
\bottomrule
\end{tabularx}
\end{center}

\noindent\textbf{注意：科目代码逐年变动。}上表的科目代码是\textbf{从各年原卷封面上逐字读出来的}，
2014 年为 818，2015 年为 819，2016 年为 817，2017/2018 年为 819，2019 年为 821，
2020 年起才稳定为 818。**不要因为"现在考 818"就把早年试卷的代码改掉**，
章标题一律按当年代码如实标注。

\noindent\textbf{重要提示：}
\begin{itemize}
  \item 早年试卷（2014--2015）含明渠流、流体机械等\textbf{现已不在初试范围}的内容，
        做这些题时\textbf{跳过超纲题}（正文已逐题标注"超纲"）。
  \item \textbf{2019 年以后并不是"一套固定模板"}：2019、2020 两年另有\textbf{作图题与证明题}，
        2021 年起才稳定为"选择 + 判断 + 填空 + 简答 + 计算"五大块。
        但共同点是\textbf{计算题分值最高}（50--60 分），选择、判断、填空合计约 60--80 分。
  \item \textbf{官方答案只有 2014--2020 年。}2021、2022 年以及 2026 年回忆版\textbf{原资料未附答案}，
        这几年的参考答案由编者按流体力学标准解法独立给出，并在该章章首明确标注，仅供自检使用。
  \item 2026 年真题为考生回忆版，题干文字可能与原卷有出入，但\textbf{考点与题型具有高度参考价值}。
\end{itemize}

\section*{三、命题规律（来自考情分析，务必记住）}

\begin{enumerate}
  \item \textbf{真题重复率很高}：同一考点会换数字、换题型反复考。
  \item \textbf{真题大部分源于课后习题的改编}——课后习题必须逐题过关（见《课后习题》分册）。
  \item 计算题集中在\textbf{伯努利方程、动量方程、管路水力计算}三大块（详见《教材精讲》第 4、5 章）。
  \item 重视基础，新题型很少；题量适中、计算量适中、总体难度中等。
\end{enumerate}

\section*{四、答案说明}

参考答案统一置于书末"参考答案"部分，按年份排列。计算题答案给出\textbf{公式 $\to$ 代入 $\to$ 结果 $\to$ 单位}的完整过程；
选择题、判断题给出正确选项与\textbf{一句话理由}。凡属旧大纲超出 818 范围的题目，答案中均标注"（旧大纲，可跳过）"。
